\section{Beyond presheaves: sites, subtopoi and theories}\label{sec:sites}

This section is lighter: it explains where the other decks of Caramello's programme come from, and states the fifth bridge. It can be skipped on a first reading, except for \S\ref{ssec:alexandrov}, which \S\ref{sec:preorders} uses.

\subsection{Sheaves on a space}

Let $X$ be a topological space. A \emph{sheaf} on $X$ assigns a set $F(U)$ to every open set $U$, with restriction maps $F(U)\to F(V)$ for $V\subseteq U$, such that elements can be \emph{glued}: if $U = \bigcup_i U_i$ and $x_i\in F(U_i)$ agree on the overlaps $U_i\cap U_j$, there is exactly one $x\in F(U)$ restricting to each $x_i$. Continuous functions, sections of a covering space, solutions of a differential equation give sheaves. The sheaves on $X$ form a topos $\Sh(X)$.

A sheaf is a presheaf on the preorder of open sets (with arrows $V\to U$ for $V\subseteq U$), subject to the gluing condition. In general $\Sh(X)$ is \emph{not} a presheaf topos, and its negation is computed differently: a sub-sheaf of the terminal sheaf is an open set $U$, and $\neg U$ is the interior of its complement. Then $\Sh(X)$ is Boolean iff every open set is closed, and De Morgan iff $X$ is \emph{extremally disconnected}, i.e.\ the closure of every open set is open \cite{johnstone-demorgan}. This is a third arch, which we use in \S\ref{sec:preorders}.

\subsection{Alexandrov spaces}\label{ssec:alexandrov}

Let $(P,\le)$ be a preorder. The \emph{Alexandrov topology} on $P$ has as open sets the \emph{up-sets}: the subsets $U$ such that $p\in U$ and $p\le q$ imply $q\in U$. Arbitrary intersections of open sets are open, every point $p$ has a smallest open neighbourhood $\uparrow p = \{q : p\le q\}$, the closed sets are the down-sets, and the closure of a set $S$ is $\downarrow S = \{p : p\le s\text{ for some }s\in S\}$. We write $X_P$ for this space.

\begin{proposition}\label{prop:alexandrov}
$\Sh(X_P)\simeq[P,\Set]$, the category of covariant functors $P\to\Set$, i.e.\ $\PSh(P^{\op})$.
\end{proposition}

\begin{proof}[Sketch]
A sheaf $F$ is determined by its values on the basic opens $\uparrow p$, since every open set is the union of the $\uparrow p$ it contains and gluing is unique. If $p\le q$ then $\uparrow q\subseteq\uparrow p$, so restriction gives a map $F(\uparrow p)\to F(\uparrow q)$: a functor $P\to\Set$. Conversely a functor $G : P\to\Set$ gives the sheaf $F(U) = \lim_{p\in U}G(p)$ (compatible families). The two constructions are inverse to each other.
\end{proof}

So Alexandrov spaces give decks: the same topos is presented by a space ($X_P$) and by a small category ($P^{\op}$).

\subsection{Grothendieck topologies}

Grothendieck's insight was that the gluing condition only needs to know which families of arrows \q{cover} an object. A \emph{Grothendieck topology} $J$ on a small category $C$ chooses, for each object $U$, a set $J(U)$ of \emph{covering sieves} on $U$, subject to three axioms (the maximal sieve covers; covering sieves are stable under pulling back along arrows; a sieve that is covering locally on a covering sieve is covering) \cite[slide 34]{caramello-mines}. A \emph{sheaf} for $J$ is a presheaf in which compatible families on a covering sieve glue uniquely, and the sheaves form the topos $\Sh(C,J)$. The pair $(C,J)$ is a \emph{site}. With the \q{trivial} topology, in which only maximal sieves cover, every presheaf is a sheaf and $\Sh(C,J) = \PSh(C)$; with $C$ the open sets of $X$ and covering sieves the families covering $U$ in the usual sense, $\Sh(C,J) = \Sh(X)$.

Every Grothendieck topos has many sites: this is the source of decks.

\subsection{Subtopoi and finer topologies}

A \emph{subtopos} of $\mathcal E$ is a full subcategory closed under limits, whose inclusion has a left adjoint that preserves finite limits. For instance, for an open subset $V\subseteq X$, the sheaves on $V$ form a subtopos of $\Sh(X)$; and $\Sh(C,J)$ is a subtopos of $\PSh(C)$, the left adjoint being sheafification.

\begin{bridge}[Subtopoi]\label{br:sub}
The subtopoi of $\Sh(C,J)$ correspond, reversing the order, to the Grothendieck topologies on $C$ that contain $J$. An equivalence $\Sh(C,J)\simeq\Sh(D,K)$ therefore gives an order isomorphism between the topologies on $C$ finer than $J$ and those on $D$ finer than $K$.
\end{bridge}

\leanbox{\lean{subtoposOrderIso}, \lean{exists_orderIso_ici_of_sheafEquiv}.}

The second sentence is a bridge whose invariant is not a property but a whole lattice (the lattice of subtopoi), read on two sites.

\subsection{Classifying topoi}

A \emph{geometric theory} is a first-order theory whose axioms have the form $\forall\vec x\,(\varphi\Rightarrow\psi)$ with $\varphi,\psi$ built from equations and relations by $\wedge$, $\exists$ and (possibly infinite) $\vee$. Groups, rings, local rings, fields (written suitably), linear orders, torsors under a group are geometric theories. A model of a theory can be taken not only in $\Set$ but in any topos. A theory $\mathbb T$ has a \emph{classifying topos} $\Set[\mathbb T]$: a topos containing a \q{generic} model, such that the models of $\mathbb T$ in any topos $\mathcal F$ correspond to the geometric morphisms $\mathcal F\to\Set[\mathbb T]$. In particular the models of $\mathbb T$ in $\Set$ are the points of $\Set[\mathbb T]$.

\begin{example}\label{ex:classifying}
\begin{enumerate}
\item The theory with one sort and no axioms (\q{the theory of an object}) is classified by $[\mathsf{Fin},\Set] = \PSh(\mathsf{Fin}^{\op})$.
\item For a group $G$, the theory of $G$-torsors (a non-empty set with a free and transitive left action of $G$) is classified by $\PSh(\B G)$.
\item More generally, a \emph{Horn} (or cartesian) theory, such as the theory of groups, of rings, of modules over a fixed ring, of lattices, is classified by $[\fp\mathbb T\text{-mod},\Set]$, the functors on its finitely presented models \cite[ch.~6]{caramello-tst}. Such theories are said to be \emph{of presheaf type}.
\end{enumerate}
\end{example}

\leanbox{The library proves that every geometric theory has a classifying site (\lean{ClassifyingTopoi.exists_isClassifyingSite}), that Horn theories are of presheaf type (\lean{ClassifyingTopoi.isOfPresheafType_of_isHorn}), and the torsor example (\lean{ClassifyingTopoi.isClassifyingSite_genericTorsor}).}

Two theories with equivalent classifying topoi are called \emph{Morita equivalent}: they have \q{the same models} in every topos, even if their languages are different. This gives the richest decks of all, and Caramello's programme is to read invariants of the common classifying topos on both theories. A \emph{quotient} of $\mathbb T$ is a theory obtained by adding axioms; Bridge~\ref{br:sub} becomes, through classifying topoi, a correspondence between the quotients of $\mathbb T$ and the topologies finer than the one of a classifying site \cite[ch.~3]{caramello-tst} (\lean{ClassifyingTopoi.IsClassifyingSite.exists_quotient_theory}). This is how Caramello computes, for instance, an explicit axiomatisation of the common consequences of the theory of local rings and of the theory of integral domains \cite[\S18]{caramello-lattices}: the meet of two quotients is computed as the meet of two Grothendieck topologies.
